\chapter{MỞ ĐẦU}

\section{Đặt vấn đề}

Phân tích cảm xúc (Sentiment Analysis) đã trở thành một lĩnh vực ngày càng quan trọng trong xử lý ngôn ngữ tự nhiên (NLP), cho phép hiểu tự động các ý kiến và cảm xúc được thể hiện trong văn bản. Với sự phát triển nhanh chóng của mạng xã hội và các nền tảng trực tuyến, hàng triệu bình luận được tạo ra mỗi ngày, chứa đựng những thông tin quý giá về ý kiến, sở thích và cảm xúc của người dùng.

\section{Lý do chọn đề tài}

Mặc dù có nhiều tiến bộ trong NLP, việc phân tích cảm xúc trong bình luận vẫn còn nhiều thách thức:
\begin{itemize}
    \item Ngôn ngữ không chính xác và tiếng lóng
    \item Phát hiện mỉa mai và châm biếm
    \item Ý nghĩa phụ thuộc vào ngữ cảnh
    \item Nội dung đa ngôn ngữ
\end{itemize}

\section{Mục tiêu nghiên cứu}

Mục tiêu chính của nghiên cứu này là:
\begin{enumerate}
    \item Xây dựng mô hình phân tích cảm xúc hiệu quả cho dữ liệu bình luận
    \item Xử lý các thách thức đặc thù của văn bản không chính thức
    \item Đánh giá hiệu suất mô hình trên các tập dữ liệu thực tế
\end{enumerate}

\section{Đóng góp của đề tài}

Các đóng góp chính của công việc này bao gồm:
\begin{itemize}
    \item Cách tiếp cận mới cho phân loại cảm xúc
    \item Đánh giá toàn diện trên nhiều tập dữ liệu
    \item Phân tích các trường hợp lỗi và hạn chế
\end{itemize}

\section{Cấu trúc đồ án}

Nội dung đồ án được tổ chức như sau. Chương~\ref{chap:related} trình bày các công trình nghiên cứu liên quan. Chương~\ref{chap:methodology} đề xuất phương pháp. Chương~\ref{chap:experiments} mô tả thiết lập thử nghiệm và kết quả. Cuối cùng, Chương~\ref{chap:conclusion} kết luận đồ án và thảo luận hướng phát triển.
